%% kapa: exam=1 round=35 year=2024 period=2 form=A subject=appraisal_law subject_ko=감정평가관계법규 seq=18 number=18
%% source: Q-Net 감정평가사 기출문제 (cst00309) artlSeq=5213229 file=raw/1008/5213229/2교시 문제지 원본.pdf
%% conversion: kordoc 4.24.51 -> tools/gen1.py (auto); review=proofread-sonnet; flags=none
\begin{question}{18}
감정평가 및 감정평가사에 관한 법령상 감정평가법인에 관한 설명으로 옳지 않은 것은?
\begin{choices}
  \item 감정평가법인은 전체 사원 또는 이사의 100분의 90 이상을 감정평가사로 두어야 한다.
  \item 국토교통부장관은 감정평가법인등이 장부 등의 검사를 거부 또는 방해한 경우에는 그 설립인가를 취소할 수 있다.
  \item 감정평가법인등은 토지등의 매매업을 직접 하여서는 아니 된다.
  \item 감정평가법인의 자본금은 2억원 이상이어야 한다.
  \item 감정평가법인의 대표사원 또는 대표이사는 감정평가사여야 한다.
\end{choices}
\end{question}
\answer{2}
